%%%%%%%%%%%%%%%%%%%%%%% file template.tex %%%%%%%%%%%%%%%%%%%%%%%%%
%
% This is a general template file for the LaTeX package SVJour3
% for Springer journals.          Springer Heidelberg 2010/09/16
%
% Copy it to a new file with a new name and use it as the basis
% for your article. Delete % signs as needed.
%
% This template includes a few options for different layouts and
% content for various journals. Please consult a previous issue of
% your journal as needed.
%
%%%%%%%%%%%%%%%%%%%%%%%%%%%%%%%%%%%%%%%%%%%%%%%%%%%%%%%%%%%%%%%%%%%
%
% First comes an example EPS file -- just ignore it and
% proceed on the \documentclass line
% your LaTeX will extract the file if required
\begin{filecontents*}{example.eps}
%!PS-Adobe-3.0 EPSF-3.0
%%BoundingBox: 19 19 221 221
%%CreationDate: Mon Sep 29 1997
%%Creator: programmed by hand (JK)
%%EndComments
gsave
newpath
  20 20 moveto
  20 220 lineto
  220 220 lineto
  220 20 lineto
closepath
2 setlinewidth
gsave
  .4 setgray fill
grestore
stroke
grestore
\end{filecontents*}
%
\RequirePackage{fix-cm}
%
%\documentclass{svjour3}                     % onecolumn (standard format)
%\documentclass[smallcondensed]{svjour3}     % onecolumn (ditto)
\documentclass[smallextended]{svjour3}       % onecolumn (second format)
%\documentclass[twocolumn]{svjour3}          % twocolumn
%
\smartqed  % flush right qed marks, e.g. at end of proof
%
\usepackage{graphicx}
%
% \usepackage{mathptmx}      % use Times fonts if available on your TeX system
%
% insert here the call for the packages your document requires
%\usepackage{latexsym}
% etc.
%
% please place your own definitions here and don't use \def but
% \newcommand{}{}
%
% Insert the name of "your journal" with
% \journalname{myjournal}
%
\begin{document}

\title{A game-theory modeling approach to fitness and trust dynamics in biomedical research social networks}%\thanks{Grants or other notes
%about the article that should go on the front page should be
%placed here. General acknowledgments should be placed at the end of the article.}

%\subtitle{Do you have a subtitle?\\ If so, write it here}

%\titlerunning{Short form of title}        % if too long for running head

\author{J. Mario Siqueiros-Garc\'ia         \and      Rodrigo Garc\'ia-Herrera \and         Enrique Hern\'andez-Lemus \and
        Sergio Alcal\'a  %etc.
}

%\authorrunning{Short form of author list} % if too long for running head

\institute{J. Mario Siqueiros-Garc\'ia \at
  IIMAS-UNAM. Circuito Escolar 3000, Cd Universitaria, Coyoacán, Ciudad de México, D.F. \\
              Tel.: +123-45-678910\\
              Fax: +123-45-678910\\
              \email{jmario.siqueiros@iimas.unam.mx}           %  \\
%             \emph{Present address:} of F. Author  %  if needed
           \and
           S. Author \at
              second address
}

\date{Received: date / Accepted: date}
% The correct dates will be entered by the editor


\maketitle

\begin{abstract}
\textbf{Background:} In an ideal world, access to research resources should be fair and equitative according to the proposals relevance and the researcher's academic record. We know that this is not necessarily so, specially in places where access to some resources, e.g., biological samples, is not regulated. Other factors may come into play like social connections, political power or prestige. In this work we explore the distribution of fitness and trust in a biomedical researchers collaboration network when playing a variation of an iterative prisoner's dilemma in which agents are compromised in either defecting and increasing their individual fitness or cooperating and increase mutual fitness with their neighbors. \textbf{Methods:} Fitness is a property of the each agent and trust is a property of the link between two agents. According to a pay-off matrix and a mutual trust $A_{ij}$ matrix, we get a measure of na\"{i}vit\'e or confidence for each node. If the agents' confidence is below certain value then the agent will act suspiciously and will defect, othewise it will cooperate. We tested our simulation on an Erd\"{o}s-Reny\'i, a Watts-Strogatz small-world and Barab\'asi-Albert topologies, as well as on a real biomedical research network. Agents behavior is updated in a synchronous manner. \textbf{Results:} All networks find a point of equilibrium before the $50^{th}$ iteration. Different topologies display different fitness and trust distribution. Fitness in an Erd\"{o}s-Reny\'i netwok follows a normal distribution and trust is bimodal. For a Watts-Strogatz, small world networks, both fitness and trust distributions are strongly skewed to the right. Barab\'asi-Albert topology has a heavy left-skewed distribution (resembling to a power-law) and trust is bimodal. The biomedical researchers network has fitness distribution as in a Barab\'asi-Albert, but trust is distributed as in a Watts-Strogatz, small world topology. \textbf{Discussion:} 1) The distribution of fitness in the researchers network suggests that there are mechanisms governing the network that produces an asymetric access to resources. 2) Nevetheless, trust variables behaves as in the small-world model which might reflect some sort of subordination among researchers by which they are obliged to cooperate. 3) The range in the threshold that regulates the stringency to cooperate or defect according to the agent's na\"{i}vit\'e is small, suggesting that there is a region in which a phase transition takes place from a population full of defectors to a population of cooperators 4) Finally, we would like to propose that this sort of work may help to make visible the need to develop science policies to promote a better, small world-like, fair fitness distribution. 

%% Insert your abstract here. Include keywords, PACS and mathematical
%% subject classification numbers as needed.
\keywords{Game theory \and Trust \and Models \and Biomedical research community \and Social Networks}
% \PACS{PACS code1 \and PACS code2 \and more}
% \subclass{MSC code1 \and MSC code2 \and more}
\end{abstract}

\section{Introduction}
\label{intro}

Your text comes here. Separate text sections with



\section{Biomedical research: CONACyT and FOSISS}
\label{sec:1}
CONACyT (Council of Science and Technology) is the Mexican government entity in charge of promoting the development of science and
technology. It is also in charge of creating the country's scientific
policies. In a way, CONACyT is equivalent to the National Science Foundation of the United States.\\

Among CONACyT's functions it is to develop science and technology
policies according to national needs and demands, to offer
assessorship to the different instances of the government on
scientific and technological topics, to promote the creation of
research networks among the scientific community, grant scholarships
for masters and doctoral studies, and it also manages different trusts intended to fund individuals
and groups for scientific and technological research. \\

In the year 2002 CONACyT, along with other government agencies and entities have created sectorial funds or \emph{Fondos Sectoriales}. Sectorial funds are trusts for scientific and technological research. The objective of such funds is to cover and equally promote the research capacities of different areas such as energy, agriculture, or technological innovation by means of the generation of human resources and helping research groups to consolidate. It is expected that the knowledge generated under the sponsorship of \emph{Fondos Sectoriales} to be the product of applied research that attends national public needs, and promotes economic growth.\\

\emph{FOSISS} or Sectorial Fund for Health and Social Security Research (\emph{Fondo Sectorial en Investigaci\'on en Salud y Seguridad Social}) is among such funds. FOSISS is constituted by CONACyT, SSA, IMSS and ISSSTE \footnote{SSA is the acronym for Secretariat of Health \emph{Secretar\'ia de Salud}; IMSS is the acronym for Social Security Mexican Institute (\emph{Instituto Mexicano del Seguro Social}); ISSSTE stands for Institute for Social Security and Services for State Workers (\emph{Instituto de Seguridad y Servicios Sociales de los Trabajadores del Estado})}, being all of them the major public health providers and research institutions in the country. Every year CONACyT opens a call for funds constrained to a set the health areas previously defined by a group of experts. Elegibility is open to the public and private health research sectors, nevertheless, most applicants are public universities and research institutions. 

\section{Methodology}
\label{sec:2}

\subsection{General description of what we did}

\subsection{Our fitness, trust and naivite model}

\subsection{Implementation of the model in different topologies}

\subsubsection{Erd\"{o}s-R\'enyi}

\subsubsection{Small-World}

\subsubsection{Barbási-Albert}

\subsubsection{Biomedical research community network}


\section{Results}
\label{sec:3}



\section{Discussion}
\label{sec:4}


%% Text with citations \cite{RefB} and \cite{RefJ}.
%% \subsection{Subsection title}
%% \label{sec:2}
%% as required. Don't forget to give each section
%% and subsection a unique label (see Sect.~\ref{sec:1}).
%% \paragraph{Paragraph headings} Use paragraph headings as needed.
%% \begin{equation}
%% a^2+b^2=c^2
%% \end{equation}

% For one-column wide figures use
%%\begin{figure}
% Use the relevant command to insert your figure file.
% For example, with the graphicx package use
%%  \includegraphics{example.eps}
% figure caption is below the figure
%%\caption{Please write your figure caption here}
%%\label{fig:1}       % Give a unique label
%\end{figure}
%
% For two-column wide figures use
%\begin{figure*}
% Use the relevant command to insert your figure file.
% For example, with the graphicx package use
%  \includegraphics[width=0.75\textwidth]{example.eps}
% figure caption is below the figure
%\caption{Please write your figure caption here}
%\label{fig:2}       % Give a unique label
%\end{figure*}
%
% For tables use
%\begin{table}
% table caption is above the table
%\caption{Please write your table caption here}
%\label{tab:1}       % Give a unique label
% For LaTeX tables use
%\begin{tabular}{lll}
%\hline\noalign{\smallskip}
%first & second & third  \\
%\noalign{\smallskip}\hline\noalign{\smallskip}
%number & number & number \\
%number & number & number \\
%\noalign{\smallskip}\hline
%\end{tabular}
%\end{table}


%\begin{acknowledgements}
%If you'd like to thank anyone, place your comments here
%and remove the percent signs.
%\end{acknowledgements}

% BibTeX users please use one of
%\bibliographystyle{spbasic}      % basic style, author-year citations
%\bibliographystyle{spmpsci}      % mathematics and physical sciences
%\bibliographystyle{spphys}       % APS-like style for physics
%\bibliography{}   % name your BibTeX data base

% Non-BibTeX users please use
\begin{thebibliography}{}
%
% and use \bibitem to create references. Consult the Instructions
% for authors for reference list style.
%
\bibitem{RefJ}
% Format for Journal Reference
Author, Article title, Journal, Volume, page numbers (year)
% Format for books
\bibitem{RefB}
Author, Book title, page numbers. Publisher, place (year)
% etc
\end{thebibliography}

\end{document}
% end of file template.tex

